\section{Related Work}
\label{sec:related}

\paragraph{Entity selection for knowledge-grounded LLMs.}
Knowledge graphs support LLM-based question answering and agents by providing relational context for retrieval and reasoning~\cite{luo2024rog,ghandi2026graphwalk}. A query may depend on intermediate entities connecting its starting point to its answers, making the selection of supporting context a problem beyond identifying plausible endpoints~\cite{luo2024rog}. STaRK evaluates entity retrieval under textual and relational query constraints, while G-Retriever constructs a connected graph context through prize-collecting Steiner tree optimization~\cite{wu2024stark,he2024gretriever}. Our study examines entity selection within a fixed candidate pool containing every reference entity. The target is the union of entities on all routes permitted by the query, including intermediate and branching entities. Evaluating this set directly allows us to characterize the completeness and selectivity of the relational context retained by the model.

\paragraph{Graph structure and task comprehension.}
Graph-reasoning benchmarks provide context for evaluating how this selection behavior varies across structural configurations~\cite{fatemi2024talklikegraph,dai2025graphpatterns,xu2026graphomni}. Talk like a Graph studies graph encodings, task types, and graph structure; Xinnan Dai et al. examine pattern comprehension and discovery from terminological and topological descriptions~\cite{fatemi2024talklikegraph,dai2025graphpatterns}. GraphOmni jointly varies graph families, serialization formats, and prompting schemes across graph-theoretic tasks~\cite{xu2026graphomni}. We use directed motifs and reference sizes to organize entity-selection profiles, interpreting structural differences alongside query semantics and the evidence visible to the model.

\paragraph{Availability and representation of relational evidence.}
The graph information presented to an LLM is another central factor in KG-based reasoning~\cite{dai2024kgunderstanding,tian2025alignment}. Xinbang Dai et al. compare KG input formats and triple organization and evaluate robustness to noisy or incomplete subgraphs using answer-based evaluation~\cite{dai2024kgunderstanding}. Tian et al. distinguish graph transformation from linearization and examine formats, orders, and templates in KG-augmented generation~\cite{tian2025alignment}. Our comparisons distinguish relation availability from the presentation of designated matched facts: Raw versus Raw+E adds explicit relations to native entity cards, while controlled triples, templated prose, and alternative card hosting vary their presentation. This distinction lets us examine entity recovery under different evidence conditions without assuming that native descriptions contain every graph relation.

\paragraph{Selection errors, distractors, and evaluation targets.}
GraphWalk reports incomplete exploration when graph context is supplied directly to the model~\cite{ghandi2026graphwalk}. TopoRAG addresses noise from heuristic neighbor expansion through retrieval constrained by induced-subgraph diameter, while Amiraz et al. quantify the distracting effect of irrelevant passages on answer-generating models~\cite{wu2026toporag,amiraz2025distracting}. Our disconnected/adjacent-context comparison examines query-irrelevant neighboring entities while preserving reference availability within a fixed-size candidate pool. This provides a structural distractor condition; adjacency alone is not a definition of irrelevance. Li and Ouyang show that the relationship between knowledge-selection quality and generation quality depends on the generator and task~\cite{li2025selection}. We therefore evaluate route coverage, excess selection, and complete-set recovery directly through recall, precision, and exact recovery, keeping these outcomes distinct from minimum sufficient evidence and final-answer correctness.

\paragraph{Relationship to structural tool selection.}
Tool selection provides a related perspective on retaining elements that compose to support a request~\cite{anonymousToolGraph}. Tool Graph Complexity distinguishes reference-set non-minimality from dependency composition and analyzes motif-associated selection difficulty~\cite{anonymousToolGraph}. Its Appendix~O directly overlaps the KG setting: it evaluates motif and size variation with explicit relations and adjacent decoys, reporting easier KG chains, difficult incoming trees, and decoy-related inclusion errors~\cite{anonymousToolGraph}. We acknowledge these prior KG findings and center the present analysis on entity-route recovery across relation availability, matched-fact presentation, and distractor conditions. Protocol differences and unresolved task/response overlap limit direct cross-study score comparisons; the analyses do not establish a common internal failure mechanism for tool and entity selection.
